\section{Conclusion}
\label{sec:conclusion}

The tested state representations do not preserve the task success of full
interaction history when used as its replacement. Retaining recent records
raises the solved rate in both model-written and hybrid state policies,
while the tested inline protocol and task-adapted schema do not recover
the full-history result. These findings identify useful information in the
retained interaction record without establishing the incremental benefit
of state over an equally sized history window.

Resource savings depend on the objective. Smaller peak prompts coexist
with higher conditional successful-execution costs, whereas all-execution
cost per success can improve when recent history is retained. State-based
context management should therefore be evaluated through task completion,
token categories, and cost together. The next discriminating comparison is
a matched history-only policy, followed by replication across models and
workloads before broader claims about structured memory are warranted.
